\documentclass{article}
\usepackage{imakeidx}  % More flexible than makeidx
\makeindex
\usepackage{hyperref}  % Optional: for clickable links
\usepackage{amsmath}        % Adds mathematical typesetting support
\usepackage{amssymb}
\usepackage{amsthm}         % Required for theorem environments
\usepackage{graphicx}
\usepackage[margin=1.5cm]{geometry}
\usepackage{siunitx}
\usepackage{tikz}

\title{Luna math}
\author{Filip Jesenšek}
\date{\today}

\begin{document}
\maketitle

\newpage

\section{Problem lociranja kraterjev}

\subsection{Matematičen opis problema lociranja točk}
Množico točk na enotski n-sferi bom poimenoval $S^n$. Definicija te množoce bi bila naslednja:

\begin{equation}
	S^n = \{x \in \mathbb{R}^{n+1} : \, ||x|| = 1\}
	\label{eq:definition_sphear}
\end{equation}

Tukaj se bomo seveda ukvarjali samo z množico točk na 2-sferi $S^2$

Išemo način kako neko točko $T$ z koordinatami $(x, y)$ na pridobljeni sliki, preslikamo na sfero $S^2$, 
ki ima neko preferenčno orientacijo (sever, jug, vzhod, zahod). Tukaj smo seveda že takoj predpostavili,
da je Lunina površina sfera, oziroma ekvivalentno, da je njena ecentričnost $e = 0$, 
v resnici je $e \approx 0.05$ (bolj okrogla skoraj ne more biti).

\subsection{Čez objektiv}
Preden se lahko sploh pogovarjamo o rotacijah različnih koordinatnih sistemov, je treba pogledati kako sama 
leča v teleskopu (ki slika luno) "deformira" sliko. To lahko preprosto modeliramo z linearno transformacijo,
ki raztegne osi in aplicira strig. V splošnem lahko tako matriko zapišemo kot.

\begin{equation}
	M = 
	\begin{pmatrix}
		f_x & s \\
		0 & f_y
	\end{pmatrix}
	\label{eq:lense_deformation}
\end{equation}

To mogoče nebo treba uporabiti, saj so lahko leče tako "dobre" da je $M = I$ ($M$ kar identiteta).

\subsection{Iz ne deformirane slike na sfero}
Tukaj je naša naloga, da našo 2D sliko spravimo na sfero. Za enkrat se ne ukvarjamo z tem, kateri del sfere smo
slikali, to nas čaka v naslednjem razdelku. Problem lahko rešimo na način, da najprej pogledamo, kaj se
zgodi z vektroji v 3D prostoru, ko jih slikamo z kamero. Ko to ugotovimo, se lahko vprašamo kaj bi lahko bila
obratna transformacija.

Že sama intuicija pove, da samo dejanje slikanja je v resnici neke vrste projekcija. V tem primeru smo
množico točk iz $S^2$ preslikali na ravnino $XY$. Izbira ravnine je seveda arbitrarna, ni nobenega posebnega
razloga zakaj bi si izbrali ravnino $XY$. To lahko formalo zapišemo, na nasledni način. Naj bo
$v \in S^2$ in naj bo $P$ projektor na $XY$ ravnino. 

\begin{equation}
	P v = w
	\label{eq:projection}
\end{equation}

Projektor v standardni bazi lahko trivialno zapišemo kar eksplicitno.

\begin{equation}
	P = 
	\begin{pmatrix}
		1 & 0 & 0 \\
		0 & 1 & 0 \\
		0 & 0 & 0
	\end{pmatrix}
	\label{eq:projector}
\end{equation}

Vektor $w$ si lahko predstavljamo, kot vektor ki ga vidimo na naši sliki. To ni več 3D objekt v
prostoru, ampak živi samo v naši 2D sliki. Mi imamo $w$ podan 
(to je una stvar, ki jo merimo/slikamo), sepravi naloga je izračunati $v$. Če nalogo bolje vbesedimo, sprašujemo so po takih vektorjih $v \in S^2$, ki po projekciji $P$ postanejo $w$ v 
ravnini $XY$. Tukaj bi vsak z veseljem množili $P^{-1}$ z leve obi strani, dobil bi
$v = P^{-1} w$. Vendar, če se malo vstavimo in bolje analiziramo ugotovimo, da $P^{-1}$
sploh ne obstaja, saj je $\det P = 0$. To pomeni, da je matrika $P$ neobrnjiva. Tako da na žalost bomo morali
iti po komponentah (ker je matrika tako preprosta bo to trivialno). Če napišemo (\ref{eq:projection}) eksplicitno
dobimo naslednje.

\begin{equation}
	\begin{pmatrix}
		1 & 0 & 0 \\
		0 & 1 & 0 \\
		0 & 0 & 0
	\end{pmatrix}
	\begin{pmatrix}
		\sin \vartheta \cos \varphi \\
		\sin \vartheta \sin \varphi \\
		\cos \vartheta
	\end{pmatrix}
	=
	\begin{pmatrix}
		x \\
		y \\
		z	
	\end{pmatrix}
	\label{eq:projector_components}
\end{equation}

Znano je da lahko vsak element iz $S^2$ napišemo na ta način z polarnimi koti. 
Tukaj $z$ komponenta seveda odpade, saj na sliki "ne vidimo globine". Ko se znebimo nepotrebne $z$ komponente
na ostane samo 2 krat 2 identiteta. Kot obljubljeno projekcija deluje trivialno na $v$ in
dobimo naslednji sistem enačb.

\begin{equation}
	\begin{aligned}
		x = \sin \vartheta \cos \varphi \\
		y = \sin \vartheta \sin \varphi 
	\end{aligned}
	\label{eq:system_of_equations}
\end{equation}

\newpage

Z deljenjem enačb in malo obračanja dobimo rešitev.

\begin{equation}
	\tan \varphi = \frac{y}{x}
	\label{eq:phi}
\end{equation}

\begin{equation}
	\sin \vartheta = \frac{x}{\cos \varphi} = \frac{y}{\sin \varphi}
	\label{eq:theta}
\end{equation}

Ko dobimo par koordinat $(x,y)$ lahko naprej po (\ref{eq:phi}) dobim $\varphi$ in nato po (\ref{eq:theta}) dobimo
še $\vartheta$.

\subsection{Pravilno orientiranje sfere}
Zdaj pa pridemo do najtežjega dela tega problema. Naša sfera ima v 3D prostoru neko odlikovano os (tej osi 
rečemo sever). Naša naloga je v resnici iskanje zapisa točk (ki imajo koordiante v našem sistemu) v 
odlikovanem koordinatem sistemu. Naš problem zdaj je izključno omejen na $S^2$. Tako da da se nebomo medli,
kateri sistem je kateri, bom pisal $v$ točke, ki so v našem sistemu in $u$ točke v odlikovanem sistemu. 
Da vse točke iz slike preslikamo na sfero v odlikovanem koordinatem sistemu, moramo izbrati 2 referenčni točki, 
ki vemo kam se preslikajo. Da enolično določimo preslikavo potrebujemo dve točki, saj je problem 2 dimenzionalen
koordinati sta $\varphi$ in $\vartheta$ naša vez (katera nam zmanjšuje dimenzijo problema) pa je radij lune.
Za rešitev tega problema bomo uporabili kvaternione. Vsi kvaternioni s katerimi se bomo ukvarjali bodo enotski.


Recimo da poznamo naslednje točke $v_{1}, v_{2}, u_{1}, u_{2}$ kjer gre $v_i \mapsto u_i$ po transformaciji.

Rotacijske kvaternione $q \in \mathbb{H}$ lahko zapišemo na naslednji način.
\begin{equation}
	q = \cos \frac{\phi}{2} + \vec{n} \sin\frac{\phi}{2}
	\label{eq:rot_quat}
\end{equation}
To je splošen zapis kvaterniona, ki če konjugiramo nek vektor z njim, dobimo rotacijo za kot $\phi$ okoli osi
$\vec{n}$ ($\vec{n}$ je seveda normiran na 1). Sepravi išemo tak $q$ da bo $q v_i q^{-1} = u_i$ za dane vektorje
$v_1, v_2$ in $u_1, u_2$. Rešitev lahko poiščemo tako, da pogledamo kateri


\end{document}
